\documentclass[10pt,a4paper]{article}

\usepackage{iftex}
\ifPDFTeX
  \usepackage[utf8]{inputenc}
  \usepackage[T1]{fontenc}
\fi
\usepackage[french]{babel}
\usepackage{lmodern}
\renewcommand{\familydefault}{\sfdefault}
\usepackage[top=1.1cm, bottom=1.1cm, left=1.4cm, right=1.4cm]{geometry}
\usepackage{titlesec}
\usepackage{enumitem}
\usepackage[hidelinks]{hyperref}
\usepackage{xcolor}

\pagestyle{empty}
\setlength{\parindent}{0pt}
\setlength{\parskip}{0pt}

\titleformat{\section}
  {\normalsize\bfseries\uppercase}
  {}
  {0em}
  {}
  [\vspace{-0.65em}\rule{\linewidth}{0.5pt}]
\titlespacing*{\section}{0pt}{0.55em}{0.35em}

\setlist[itemize]{leftmargin=1.1em, itemsep=0.12em, parsep=0pt, topsep=0.15em, label=--}

\begin{document}

% === HEADER ===
\begin{center}
  {\LARGE\bfseries Santiago FUENTES LLANOS}\\[0.25em]
  % === CV_TITLE_START ===
  {\large\bfseries Stage -- Analyste Cybersécurité \& Développement Systèmes}\\[0.2em]
  % === CV_TITLE_END ===
  % === CV_AVAILABILITY_START ===
  {\small\textbf{Disponibilité :} Du 03/05 au 16/07 (3 mois, prolongeable l'été) \textperiodcentered{} Stage conventionné (Université Paris-Saclay)}\\[0.15em]
  % === CV_AVAILABILITY_END ===
  {\small santi.fuentes.llanos@gmail.com \textperiodcentered{} +33 7 79 73 74 19 \textperiodcentered{} Bures-sur-Yvette (Mobilité : Île-de-France)}\\
  {\small \href{https://github.com/Devilijho}{github.com/Devilijho} \textperiodcentered{} \href{https://www.linkedin.com/in/santiago-f-ll}{linkedin.com/in/santiago-f-ll}}
\end{center}

\vspace{0.25em}
% === CV_ACCROCHE_START ===
{\small\itshape
Étudiant en 2\textsuperscript{ème} année de BUT Informatique (Réseaux, Cryptographie et Sécurité) avec le Tronc Commun de l'École 42 validé, je recherche un stage de 3 mois du 03/05 au 16/07 (prolongeable durant l'été). Passionné par l'ingénierie logicielle, la sécurité bas niveau et les systèmes, je maîtrise les environnements UNIX, le développement C/C++/TypeScript et l'architecture réseau. Rigoureux et autonome, je souhaite contribuer concrètement à vos projets techniques.
}
% === CV_ACCROCHE_END ===

% === SKILLS_SECTION_START ===
\section*{Compétences Techniques}
{\small
\textbf{Langages :} C, C++, Assembly (x86), Java, SQL (Oracle / SQLite), HTML5/CSS3, JavaScript, TypeScript, Notions de C\#\\
\textbf{Sécurité \& Débogage :} GDB, Valgrind (Memcheck, Helgrind), Analyse mémoire, Reverse Engineering (notions)\\
\textbf{Systèmes \& DevOps :} Linux (UNIX), Windows, MacOS, Git, Docker, VirtualBox\\
\textbf{Frameworks \& Outils :} Node.js, React, Three.js, Eclipse, Visual Paradigm, Modélisation (MCD/SR)
}
% === SKILLS_SECTION_END ===

% === PROJECTS_SECTION_START ===
\section*{Projets Informatiques}
{\small
\textbf{Webserv -- École 42} --- \textit{Serveur HTTP \& Sécurité Réseau en C++}
\begin{itemize}
  \item Développement d'un serveur HTTP/1.1 conforme RFC avec multiplexation d'I/O non-bloquantes (\texttt{select}/\texttt{poll}/\texttt{kqueue}).
  \item Sécurisation du parsing des requêtes contre les attaques par déni de service (DoS), corruption de flux et injection de payloads.
  \item Gestion rigoureuse des descripteurs de fichiers, des sessions TCP et des codes d'état HTTP.
\end{itemize}

\vspace{0.25em}
\textbf{Cub3d -- École 42} --- \textit{Moteur 3D \& Gestion Mémoire Bas Niveau en C}
\begin{itemize}
  \item Implémentation d'un moteur de rendu graphique par raycasting (algorithmique bas niveau et arithmétique de pointeurs).
  \item Audit de sécurité mémoire : prévention proactive des débordements de tampon (buffer overflows) et validation stricte avec Valgrind (zéro fuite mémoire / segmentation fault).
  \item Parsing défensif des fichiers de configuration de scènes (validation stricte des entrées utilisateurs).
\end{itemize}

\vspace{0.25em}
\textbf{Game Jam IUT Paris-Saclay} --- \textit{Hackathon de Programmation (36h)}
\begin{itemize}
  \item Conception rapide et sécurisée d'un exécutable complet sous contrainte de temps stricte.
  \item Structuration de la logique applicative, gestion des états et tests fonctionnels pour éliminer les bugs d'exécution critiques.
\end{itemize}
}
% === PROJECTS_SECTION_END ===

\section*{Formation}
{\small
\textbf{2025 -- 2028 (En cours) : BUT 2 Informatique} --- \textit{Université Paris-Saclay (IUT d'Orsay)}
\begin{itemize}
  \item Parcours B : Réseaux avancés, Cryptographie, Sécurité, administration et supervision des systèmes.
\end{itemize}
\textbf{Avril 2026 : École 42 Paris (Tronc Commun validé)}
\begin{itemize}
  \item Apprentissage par projets, peer-learning, programmation systèmes \& web avancée (C/C++/TypeScript/Docker).
\end{itemize}
\textbf{Juin 2025 : Baccalauréat Général} --- \textit{Mention « Bien », Spécialités Mathématiques et NSI}
}

\section*{Expérience Professionnelle}
{\small
\textbf{2022 -- 2023 : Responsable Opérationnel (E-commerce)} --- \textit{Mimitos Click}
\begin{itemize}
  \item Gestion de la logistique et du marketing digital pour une boutique en ligne.
  \item Suivi des commandes, résolution de problèmes et maintien de la qualité de service client.
\end{itemize}
}

\section*{Soft Skills \& Centres d'Intérêt}
{\small
\begin{itemize}
  \item \textbf{Soft Skills :} Rigueur scientifique, capacité d'adaptation, autonomie, travail en équipe (Peer-learning).
  \item \textbf{Centres d'intérêt :} Cybersécurité offensive (défis Capture The Flag / Reverse Engineering) et pratique de la guitare.
\end{itemize}
}

\section*{Langues}
{\small
\textbf{Espagnol :} Langue maternelle \quad \textperiodcentered{} \quad \textbf{Anglais :} Niveau C1 (Professionnel) \quad \textperiodcentered{} \quad \textbf{Français :} Niveau C1 (Professionnel)
}

\end{document}
